% TOP_PROBLEM: {"id": 20001028, "title": "Exact deficit identity for recursive tight-cycle constructions", "category_id": 2, "category": {"id": 2, "name": "combinatorics", "display_name": "Combinatorics", "description": "Counting problems, graph theory, discrete structures.", "slug": "combinatorics", "order_index": 2, "created_at": "2026-07-31T15:26:25.671Z"}, "status": "partially_solved", "problem_number": "AIM-COMBINATORICS-0153", "set_id": 14}
% FIRST_POSTED: 2026-10-01
% ENTRY_KIND: statement_only
% UnsolvedMath ID 20001028; source catalogue code AIM-COMBINATORICS-0153
% !TeX program = lualatex
% Definition and sources only; this file is not an LLM solution attempt.
\documentclass[11pt,a4paper]{article}
\usepackage[margin=25mm]{geometry}
\usepackage{fontspec}
\setmainfont{lmroman10-regular.otf}[BoldFont=lmroman10-bold.otf,
 ItalicFont=lmroman10-italic.otf,BoldItalicFont=lmroman10-bolditalic.otf]
\usepackage{amsmath,amssymb,mathtools}
\usepackage{unicode-math}
\setmathfont{latinmodern-math.otf}
\AtBeginDocument{\let\setminus\smallsetminus}
\usepackage{xurl,hyperref}
\hypersetup{colorlinks=true,urlcolor=blue}
\setcounter{secnumdepth}{0}
\setlength{\parindent}{0pt}
\setlength{\parskip}{5pt}
\setlength{\emergencystretch}{3em}
\begin{document}
\section*{Exact deficit identity for recursive tight-cycle constructions}\label{20001028}
{\small UnsolvedMath ID 20001028 · AIM-COMBINATORICS-0153 · Combinatorics · AIM Workshop Problem Lists}
\subsection{Definitions and mathematical statement}
The tight $3$-uniform five-cycle $C_5^{(3)}$ has vertices $1,\ldots,5$
and the five edges
\[
                 123,\quad234,\quad345,\quad451,\quad512.
\]
Each edge consists of three consecutive vertices in the cyclic order;
the notation suppresses braces. A simple triple system $H$ is
$C_5^{(3)}$-free if no injective map from these five vertices takes
all five displayed edges to edges of $H$. Additional triples among
those vertices are permitted in a copy, so the condition is not
induced freeness.
Define $\operatorname{ex}(n,C_5^{(3)})$ as the largest edge count of
such a free hypergraph on $n$ vertices and write
\[
 \pi(C_5^{(3)})=\lim_{n\to\infty}
       \frac{\operatorname{ex}(n,C_5^{(3)})}{\binom n3}.
\]
The limit exists, and the conjectured value is $2\sqrt3-3$.

Thus the upper-bound target is: for every $\varepsilon>0$, all
sufficiently large triple systems with density above
$2\sqrt3-3+\varepsilon$ contain the five-cycle. The proposed lower
density comes from recursively dividing a vertex set into parts
$A,B$, taking all triples with exactly two vertices in $A$ and one
in $B$, and applying the same construction inside $B$.
The problem asks whether the optimal limiting density is already
achieved by this type of construction. The forbidden object has all
five cyclic edges; deleting an edge would give a different extremal
problem and a different target.
\subsection{Short English statement}
Determine whether the tight five-cycle has triple-system Turán density exactly two square roots of three minus three.
\subsection{Sources}
\begin{itemize}
\item[S1] ULAM.AI and the UnsolvedMath Contributors, supplied problems.json, record 20001028 (AIM-COMBINATORICS-0153), AIM Workshop Problem Lists. Catalogue identity and source transcription; CC BY 4.0.
\url{https://huggingface.co/datasets/ulamai/UnsolvedMath}.
\item[S2] American Institute of Mathematics, Hypergraph Turan problem, item 33.1. Original workshop question underlying AIM-COMBINATORICS-0153.
\url{http://aimpl.org/hypergraphturan/3/}.
\end{itemize}
\end{document}
